% ===========================================================================
\section{The semantic kernel $\Ksem$}\label{sec:kernel}
% ===========================================================================

The kernel (\code{include/twin/kernel}, \code{src/kernel}) is the only component
that computes semantic successors. Its library target links only the IR value
types and the dependency-free core (errors, exact time). It performs no I/O,
networking, persistence or logging. This section defines the kernel transition
system as the functions the code computes.

\subsection{Executable model}

\begin{definition}[Kernel model]\label{def:kmodel}
\code{Model::create(M)} succeeds for a well-formed $M$ (it re-runs
\code{ir::validate}) and produces $\hat M$, which consists of $M$ unchanged plus, for
every atom $(i,j,\bowtie,c)$ of every invariant and guard, the \emph{tick atom}
$(i,j,\bowtie,\hat c)$ with $\hat c=c\cdot R$. The multiplication is checked: if
$|c|>2^{31}-1$ or the product overflows, creation fails. The remaining data
(outgoing-transition lists sorted by index, a label table) are derived indexes
that change no semantics.
\end{definition}

\begin{definition}[Kernel configurations]\label{def:kconf}
A kernel configuration is the C++ value
\begin{lstlisting}
struct Configuration {
    uint32_t             location;  // l
    std::vector<int64_t> clocks;    // v, in ticks; clocks[i-1] is clock i
    int64_t              time;      // theta, in ticks
};
\end{lstlisting}
Mathematically, $\kappa=(l,c,\tau)\in\Nat\times\Int^n\times\Int$. Its value of clock
$i$ is $c_{i-1}$ for $i\ge1$ and $0$ for the reference clock. The set of
\emph{admissible} kernel configurations is
\[
  \W(\hat M)=\{(l,c,\tau)\mid l<k,\ |c|=n,\ 0\le c_{i}\le\tau\le\Tmax,\ (l,c,\tau)\text{ satisfies }\hat J_l\}.
\]
This is exactly what \code{check\_configuration} accepts.
\end{definition}

\subsection{Kernel functions}

The following functions are those of \code{src/kernel/semantics.cpp}. They are
pure: they read their arguments and the immutable $\hat M$ and return a value
or an error, and they have no other effect.

\begin{definition}[Atom satisfaction]\label{def:katom}
$\kappa$ satisfies the tick atom $(i,j,\bowtie,\hat c)$ iff
$\mathit{val}_\kappa(i)-\mathit{val}_\kappa(j)\bowtie\hat c$, computed in \code{int64}
(\code{atom\_holds}). $\kappa$ satisfies a conjunction iff it satisfies every atom
(\code{satisfies}).
\end{definition}

\begin{definition}[Kernel transition functions]\label{def:kfun}\leavevmode
\begin{lstlisting}
initial_configuration(M) =
  c0 := (iota, 0^n, 0);   ok(c0) if c0 satisfies J_iota else err
delay(M, (l,c,t), d) =
  if d < 0                              -> err(InvalidArgument)
  if (l,c,t) does not satisfy J_l        -> err(InvariantViolation)
  t' := t + d; c' := c + d   (checked; any value > T_max -> err(Overflow))
  if (l,c',t') does not satisfy J_l      -> err(InvariantViolation)
  ok((l,c',t'))
fire(M, (l,c,t), j) =                    // transition t_j = (id,src,tgt,a,g,Y)
  if j >= m                              -> err(InvalidArgument)
  if l != src                            -> err(NotEnabled)
  if (l,c,t) does not satisfy g          -> err(NotEnabled)
  c' := c with c'[y-1] := 0 for y in Y
  if (tgt,c',t) does not satisfy J_tgt   -> err(NotEnabled)
  ok((tgt,c',t))
propositions(M, (l,c,t)) = { i | p_i.location == l }      // = {l}
\end{lstlisting}
\end{definition}

\begin{definition}[Kernel LTTS $\Ksem(M)$]\label{def:ksem}
$\Ksem(M)$ is the LTTS over the tick domain with states $\W(\hat M)$, initial state
the result of \code{initial\_configuration}, labels $\{0,\dots,m-1\}$ (internal
iff $\alpha_j=\tauact$), propositions $P$, labelling \code{propositions}, and
\[
  \kappa\trans{\delta}\kappa' \iff \code{delay}(\hat M,\kappa,\delta)=\ok(\kappa'),
  \qquad
  \kappa\trans{j}\kappa' \iff \code{fire}(\hat M,\kappa,j)=\ok(\kappa').
\]
The delay label $\delta\in\Nat$ is in ticks. When kernel and IR are compared,
$\delta$ corresponds to the delay $\delta/R\in\TR$ (\cref{sec:iso}).
\end{definition}

\begin{remark}[Derived kernel functions]
\code{enabled\_transitions}, \code{discrete\_successors}, \code{timed\_step}
($=\code{fire}\circ\code{delay}$), \code{enabling\_window}, \code{observe} and the
exploration functions are defined in terms of \code{delay} and \code{fire} or
characterised by them (\cref{lem:window,prop:monitor}). They add no transitions.
\end{remark}

\subsection{Arithmetic safety}

\begin{lemma}[No overflow]\label{lem:overflow}
For $\kappa\in\W(\hat M)$, every \code{int64} expression evaluated by
\code{atom\_holds}, \code{delay} and \code{fire} is computed without overflow,
or the function returns the explicit error \code{ArithmeticOverflow}.
\end{lemma}
\begin{proof}
Clock values lie in $[0,\Tmax]$ with $\Tmax<2^{62}$, so differences lie in
$(-2^{62},2^{62})$. Tick bounds satisfy $|\hat c|\le(2^{31}-1)\cdot10^9<2^{61}$.
Comparisons of such values are exact. In \code{delay}, the additions $\tau+\delta$
and $c_i+\delta$ use \code{\_\_builtin\_add\_overflow}, and results above $\Tmax$ are
rejected explicitly. \code{fire} assigns only the constant $0$.
\end{proof}
